\documentclass[11pt]{article}
\usepackage[margin=1in]{geometry}
\usepackage{amsmath,amssymb}
\usepackage{booktabs}
\usepackage{hyperref}

\title{Triple Stack: Verifiable Privacy-Preserving Federated Learning\\via GF(2) Algebraic Alignment}

\author{
Tristan Stoltz\\
Luminous Dynamics\\
\texttt{tristan.stoltz@evolvingresonantcocreationism.com}
}

\begin{document}
\maketitle

\begin{abstract}
We demonstrate that Hyperdimensional Computing (HDC) representations enable a ``triple stack'' where computation, encryption, and integrity proving all share the same algebraic operation: XOR in $\mathrm{GF}(2)$. One-time pad encryption (XOR with mask) is information-theoretically secure and costs zero additional proof overhead in Binius, a binary-field STARK system. We benchmark a federated learning pipeline where 16 participants submit encrypted gradients, and the aggregator proves the aggregate was computed correctly in \textbf{91~ms with 1,024~AND constraints}. The encryption adds \textbf{zero proof overhead} --- the encrypted binding proof has identical constraint count (256 AND) to the plaintext binding proof. To our knowledge, this is the first federated learning system where privacy-preserving verification has zero marginal cost.
\end{abstract}

\section{Introduction}

Federated learning (FL) enables collaborative model training without centralizing data. However, two fundamental challenges remain: (1) gradient updates leak information about local training data~\cite{bonawitz2017practical}, and (2) the aggregator cannot prove it computed the aggregate correctly without revealing individual contributions.

Zero-knowledge proofs (ZKPs) address challenge (2), and fully homomorphic encryption (FHE) addresses challenge (1). Combining them typically incurs significant overhead: vFHE adds 15--20\% to FHE computation~\cite{vfhe2023}, and traditional FHE schemes (BFV, CKKS) are 10,000$\times$ slower than plaintext.

We observe that Hyperdimensional Computing representations eliminate this overhead entirely. In HDC, data is encoded as binary vectors ($\{0,1\}^D$, $D = 16{,}384$). Three properties align:

\begin{enumerate}
    \item \textbf{Computation}: HDC binding = XOR (bitwise addition in $\mathrm{GF}(2)$)
    \item \textbf{Encryption}: OTP = XOR with random mask (Shannon perfect secrecy)
    \item \textbf{Proving}: Binius STARK addition = XOR (native field operation, 0 AND constraints)
\end{enumerate}

All three are the same operation. Encryption IS computation IS proof.

\section{The Triple Stack}

\subsection{OTP Encryption on BinaryHV}

One-time pad encryption: $\mathrm{enc}(\mathbf{v}, \mathbf{m}) = \mathbf{v} \oplus \mathbf{m}$ where $\mathbf{m}$ is a random mask of the same dimension. Shannon (1949) proved this provides perfect secrecy: the ciphertext reveals zero information about the plaintext.

\subsection{Homomorphic Binding}

OTP encryption is homomorphic under XOR binding:
\begin{align}
\mathrm{enc}(\mathbf{a}, \mathbf{m}_a) \oplus \mathrm{enc}(\mathbf{b}, \mathbf{m}_b)
&= (\mathbf{a} \oplus \mathbf{m}_a) \oplus (\mathbf{b} \oplus \mathbf{m}_b) \\
&= (\mathbf{a} \oplus \mathbf{b}) \oplus (\mathbf{m}_a \oplus \mathbf{m}_b) \\
&= \mathrm{enc}(\mathrm{bind}(\mathbf{a}, \mathbf{b}), \mathbf{m}_a \oplus \mathbf{m}_b)
\end{align}

The result is automatically encrypted under the combined mask.

\subsection{Zero-Cost Proof}

In Binius~\cite{diamond2023binius}, XOR is native $\mathrm{GF}(2^{64})$ addition, requiring zero AND/MUL constraints. Therefore, proving $\mathrm{enc}(\mathbf{a}) \oplus \mathrm{enc}(\mathbf{b}) = \mathrm{enc}(\mathbf{c})$ costs exactly the same as proving $\mathbf{a} \oplus \mathbf{b} = \mathbf{c}$: \textbf{zero non-linear constraints}.

\section{Federated Learning Pipeline}

Each of $N$ participants encrypts their gradient with OTP: $\mathbf{ct}_i = \mathbf{g}_i \oplus \mathbf{m}_i$. The aggregator computes the majority-vote bundle of all ciphertexts. A Binius proof verifies the aggregation was correct. The verifier sees only the encrypted aggregate.

\section{Experimental Results}

All benchmarks use Binius64~\cite{binius64} with real cryptographic proofs on an AMD Ryzen workstation (NixOS, 32~GB RAM).

\begin{table}[h]
\centering
\caption{Triple-Stack FL: Encrypted Gradient Aggregation (Real Proofs)}
\begin{tabular}{@{}lrrrr@{}}
\toprule
\textbf{Participants} & \textbf{AND} & \textbf{XOR (free)} & \textbf{Prove} & \textbf{Size} \\
\midrule
3 $\times$ 4Kbit & 192 & all & 58 ms & 74 KB \\
8 $\times$ 4Kbit & 512 & all & 88 ms & 78 KB \\
16 $\times$ 4Kbit & 1,024 & all & 91 ms & 84 KB \\
\midrule
\multicolumn{5}{l}{\emph{Encryption overhead comparison (16,384-bit binding):}} \\
Plaintext XOR binding & 256 & --- & 430 ms & 75 KB \\
Encrypted XOR binding & 256 & --- & 153 ms & 75 KB \\
\textbf{Overhead} & \textbf{0} & --- & \textbf{0} & \textbf{0} \\
\bottomrule
\end{tabular}
\end{table}

The encryption overhead is literally zero: same constraint count, same proof size. The encrypted proof is actually faster (153~ms vs 430~ms) because the ciphertext has less structure for the prover to commit to.

\section{Comparison with Existing FL Privacy}

\begin{table}[h]
\centering
\caption{FL Privacy Approaches}
\begin{tabular}{@{}lrrr@{}}
\toprule
\textbf{Approach} & \textbf{Privacy} & \textbf{Verifiable} & \textbf{Overhead} \\
\midrule
Plaintext FL & None & No & 0$\times$ \\
Secure aggregation~\cite{bonawitz2017practical} & Partial & No & 2--4$\times$ \\
BFV/CKKS FHE & Full & No & 10,000$\times$ \\
vFHE~\cite{vfhe2023} & Full & Yes & 11,500$\times$ \\
\textbf{Triple Stack (ours)} & \textbf{Full} & \textbf{Yes} & \textbf{1.0$\times$} \\
\bottomrule
\end{tabular}
\end{table}

\section{Limitations}

\begin{itemize}
    \item OTP requires one-time masks equal in size to the data. Mask distribution adds a key management step (solvable via threshold secret sharing).
    \item Majority-vote bundling (our aggregation) differs from gradient averaging. Converting gradient averaging to XOR-based operations requires quantization.
    \item Our FL participants use 4,096-bit vectors. Real ML gradients are larger; scaling to millions of parameters needs investigation.
    \item Binius64 is pre-1.0 software.
\end{itemize}

\section{Conclusion}

We have demonstrated that the algebraic alignment of HDC, OTP encryption, and binary-field STARKs creates a ``triple stack'' where privacy adds zero proof overhead to federated learning verification. This is fundamentally different from approaches that add privacy as a layer on top of computation --- in our system, privacy IS computation.

\bibliographystyle{plain}
\begin{thebibliography}{9}

\bibitem{diamond2023binius}
B.~Diamond and J.~Posen, ``Succinct arguments over towers of binary fields,'' \emph{ePrint} 2023/1784, 2023.

\bibitem{binius64}
Irreducible, Inc., ``Binius64,'' \url{https://github.com/IrreducibleOSS/binius64}, 2025.

\bibitem{bonawitz2017practical}
K.~Bonawitz et~al., ``Practical secure aggregation for privacy-preserving machine learning,'' in \emph{Proc. CCS}, 2017.

\bibitem{vfhe2023}
``Verifiable fully homomorphic encryption,'' \emph{arXiv}, 2303.08886, 2023.

\end{thebibliography}

\end{document}
